La cuestión afecta al estatuto explicativo de las constantes. Cuando una
teoría recibe un parámetro determinado experimentalmente, sus ecuaciones
establecen las consecuencias que se siguen de ese valor. Una derivación
estructural modifica esa relación: el valor pasa a formar parte de las
consecuencias de la construcción. La pregunta por la intensidad de un
acoplamiento o por una razón característica de la materia se convierte
entonces en la pregunta por las relaciones matemáticas que las determinan.
El tránsito del parámetro a la consecuencia constituye el argumento rector
de este artículo.

La determinación numérica adquiere, desde esta perspectiva, un contenido
constitutivo. La ecuación reúne operaciones cuya procedencia debe hacerse
explícita; su evaluación obtiene el valor que esas operaciones fijan. Una
vez establecidas las condiciones de existencia, unicidad y compatibilidad
correspondientes, la necesidad del resultado pertenece a la demostración.
Explicar por qué una constante adopta un valor significa reconstruir esa
necesidad dentro del dominio de la teoría. La precisión decimal muestra
una consecuencia cuantitativa de la estructura, mientras la derivación
expone la razón por la que se obtiene precisamente esa consecuencia.

El alcance de esta cuestión aumenta cuando varias determinaciones
comparten un mismo origen. La tesis propuesta por los autores vincula
las constantes con la organización que produce sus relaciones recíprocas.
La arquitectura común debe explicar tanto la individualidad de cada
coordenada como las dependencias que la enlazan con las restantes.
El orden de las construcciones adquiere por ello valor demostrativo:
permite reconocer qué datos se generan, qué operaciones los utilizan y
qué relaciones sobreviven a las sucesivas realizaciones. La diversidad
de valores se estudia como expresión de una unidad formal articulada.

En esta organización, la metrología desempeña una función posterior y
esencial. La construcción fija sus resultados antes de confrontarlos con
las magnitudes observadas; la experiencia examina su correspondencia
física y delimita el alcance del acuerdo empírico. La necesidad matemática se establece
respecto de las premisas y operaciones declaradas, y la realización
natural se examina mediante la identificación de observables y su
contrastación independiente. Ambas tareas contribuyen a una explicación
completa: la primera determina el resultado dentro del corpus; la segunda
estudia su pertinencia para describir la constitución cuantitativa del
mundo.

Esta perspectiva permite situar históricamente la investigación. Los
antecedentes considerados a continuación precisan distintas formas de
relacionar leyes, parámetros y observación. El presente trabajo adopta
como problema central la determinación interna de los valores y de la
estructura que los vincula. Sus desarrollos aritméticos, geométricos y
algebraicos reciben su unidad argumental de esa pregunta. La aspiración
ontológica reside en explicar cómo propiedades físicas aparentemente
independientes pueden corresponder a relaciones necesarias de una misma
organización matemática.
